% =====================================================================
%  2. EL ASISTENTE DE DEMOSTRACIÓN LEAN
% =====================================================================
\chapter{El asistente de demostración Lean}\label{sec:lean}

Este capítulo presenta Lean con un criterio concreto: explicar lo aprendido en el manejo de la herramienta y lo necesario para la formalizar el capítulo \ref{sec:formalizacion}, cosas equivalentes. Empezamos situando Lean entre los asistentes de demostración (sección \ref{sec:itp}); después describimos su fundamento teórico, la teoría de tipos dependientes, en la medida en que hace falta para entender qué significa que ``una demostración es un término'' (sección \ref{sec:tipos}); a continuación explicamos cómo se
escriben definiciones y teoremas (sección \ref{sec:escribir-lean}); cómo
funciona el modo táctico (sección \ref{sec:modo-tactico}); y cerramos con un
catálogo razonado de todas las tácticas que se usan en este trabajo (sección
\ref{sec:tacticas}) y con las herramientas de Mathlib de las que dependemos,
en particular la lógica clásica y el lema de Zorn (sección
\ref{sec:mathlib}).

\section{La demostración asistida por ordenador}\label{sec:itp}

Un asistente de demostración es un programa con dos componentes bien
diferenciadas. Por un lado, un lenguaje formal en el que el usuario
escribe definiciones, enunciados y demostraciones. Por otro, un
núcleo (\emph{kernel}): un verificador pequeño que
comprueba, paso a paso, que cada demostración es correcta según las reglas
de la lógica subyacente. Todo lo demás (editor, tácticas, automatización...) son capas de comodidad: pueden ayudar a construir la demostración, pero la única autoridad que la acepta es el núcleo.
Esta arquitectura, conocida como criterio de De Bruijn (AAA), es lo que hace
razonable fiarse del sistema: para dudar de un teorema verificado hay que
dudar de un programa de unos pocos miles de líneas, no de los cientos de
miles que forman el resto de la herramienta.

Los primeros sistemas de este tipo, como Automath (1967) y Mizar (1973),
demostraron que la idea era viable. Hoy los asistentes más usados en
matemáticas son Rocq, Isabelle/HOL y Lean. Lean \cite{lean4}, desarrollado inicialmente por Leonardo de Moura en Microsoft Research y mantenido actualmente por la Lean Focused Research Organisation, es a la vez un asistente de demostración y un lenguaje de programación funcional. En este trabajo usamos su cuarta versión, Lean~4, junto con su biblioteca matemática Mathlib \cite{mathlib}, un proyecto comunitario que reúne, verificadas, buena parte de las matemáticas de grado y posgrado. De Mathlib solo importaremos su infraestructura básica: la teoría elemental de
conjuntos, el modo táctico y, en un punto crucial, el lema de Zorn.

En la práctica, trabajar con Lean se parece a mantener un diálogo. El
usuario escribe en el editor y una ventana lateral, el \emph{InfoView}, muestra en cada momento el estado de la demostración: qué hipótesis hay disponibles y qué falta por demostrar.
Cada instrucción del usuario transforma ese estado, y la demostración
termina cuando no queda nada pendiente por probar. Veremos ese mecanismo con detalle en
la sección \ref{sec:modo-tactico}.

\section{El fundamento: la teoría de tipos}\label{sec:tipos}

La lógica sobre la que se construye Lean no es la teoría de conjuntos de
Zermelo--Fraenkel, sino una \emph{teoría de tipos dependientes} llamada
cálculo de construcciones con tipos inductivos. No necesitamos conocer este cálculo en
profundidad, pero sí entender tres de sus ideas, porque aparecen
constantemente en la formalización: 1) todo término tiene un tipo, 2) las
proposiciones son tipos, y 3) los tipos inductivos.

\subsection{Todo término tiene un tipo}

En teoría de tipos, los objetos básicos son los términos, y cada
término tiene asignado un tipo. La notación \lean{x : T} se lee
``\lean{x} es un término de tipo \lean{T}'' y es la afirmación fundamental
del sistema; juega un papel análogo al de la pertenencia $x \in A$ en teoría
de conjuntos, con una diferencia importante: un término tiene un único tipo,
que puede calcularse mecánicamente. Por ejemplo, \lean{(2 : Nat)} declara
que \lean{2} es un número natural, y \lean{Nat} es a su vez un término cuyo
tipo es \lean{Type}, el ``tipo de los tipos'' de datos.%
\footnote{Para evitar paradojas del estilo de la de Russell, \lean{Type} no
puede ser término de sí mismo: Lean organiza los tipos en una jerarquía
infinita \lean{Type 0}, \lean{Type 1}, \lean{Type 2}, \dots\ llamada
\emph{jerarquía de universos}. En este trabajo no hará falta subir nunca de
\lean{Type} (que abrevia \lean{Type 0}) y \lean{Prop}, así que no volveremos
a mencionar los universos.}

Los tipos de funciones se escriben con una flecha: una función de los
naturales en los naturales tiene tipo \lean{Nat → Nat}, y si
\lean{f : Nat → Nat} y \lean{n : Nat}, la aplicación se escribe por
yuxtaposición, \lean{f n}, sin paréntesis. Las funciones de varios
argumentos se representan encadenando flechas, que asocian a la derecha:
\lean{Nat → Nat → Nat} es el tipo de las funciones que reciben un natural y
devuelven una función \lean{Nat → Nat}; en la práctica, una función de dos
argumentos. Las funciones anónimas se escriben con la palabra clave
\lean{fun}: por ejemplo, \lean{fun n => n + 1} es la función que a cada
\lean{n} le asigna \lean{n + 1}. Esta notación aparecerá al definir modelos
concretos, como \lean{R := fun _ _ => False} (la relación que no relaciona
nada con nada; el guion bajo indica un argumento cuyo nombre no interesa).

La teoría es dependiente porque el tipo del resultado de una función
puede depender del valor del argumento. El ejemplo que nos importa es el
cuantificador universal: el tipo \lean{∀ x : A, P x} es el tipo de las
funciones que a cada \lean{x : A} le asignan un término de tipo \lean{P x},
donde \lean{P x} varía con \lean{x}. Esta lectura del $\forall$ como un tipo
de funciones tiene consecuencias muy prácticas que veremos enseguida.

\subsection{Las proposiciones son tipos}

Junto a \lean{Type}, Lean tiene otro universo llamado \lean{Prop}: el tipo
de las proposiciones. Un término \lean{p : Prop} es un enunciado
(``$2$ es par'', ``$\varphi$ es válida''), y la idea central del sistema,
conocida como correspondencia de Curry--Howard o \emph{proposiciones como
tipos}, es la siguiente:

\begin{center}
Una demostración de la proposición \lean{p} es un término de tipo
\lean{p}.
\end{center}

Es decir, las proposiciones se tratan exactamente igual que los demás
tipos, y demostrar \lean{p} consiste en construir un término \lean{h : p}.
Si no existe ningún término de tipo \lean{p}, la proposición no es
demostrable. Bajo esta correspondencia, las conectivas lógicas se convierten
en constructores de tipos y su comportamiento queda determinado por el de
los tipos correspondientes:

\begin{itemize}
  \item \textbf{Las que son funciones.} Su demostración es un procedimiento
        que transforma demostraciones en demostraciones, y usarla consiste
        en aplicarlo.
        \begin{itemize}
          \item La implicación \lean{p → q} es el tipo de las
                funciones de \lean{p} en \lean{q}: demostrarla es dar una
                función que convierte demostraciones de \lean{p} en
                demostraciones de \lean{q}. Por eso, si \lean{h : p → q} y
                \lean{hp : p}, entonces \lean{h hp : q}. El modus ponens es
                la aplicación de funciones.
          \item El cuantificador \lean{∀ x : A, P x} es el tipo de funciones
                dependientes descrito antes: demostrarlo es dar un
                procedimiento que a cada \lean{x} le asigna una demostración
                de \lean{P x}, y usarlo sobre un elemento concreto \lean{a}
                es aplicar la función: \lean{h a : P a}.
          \item La negación \lean{¬p} se define como
                \lean{p → False}, donde \lean{False} es la proposición sin
                demostraciones (un tipo vacío) y \lean{True} la que tiene una
                demostración trivial. Demostrar que \lean{p} es falsa es,
                pues, dar una función que convierta cualquier hipotética
                demostración de \lean{p} en una demostración de lo
                imposible. Este detalle será importante: nuestra negación
                modal $\mathord{\sim}\varphi := \varphi \limp \bot$ imita
                exactamente esta definición, y esa imitación hará que varios
                lemas semánticos sean ciertos ``por definición''.
        \end{itemize}
  \item \textbf{Las que son pares (o sumas).} Su demostración se construye
        empaquetando piezas, y usarla consiste en extraerlas.
        \begin{itemize}
          \item La conjunción \lean{p ∧ q} es el tipo de los pares: una
                demostración suya es un par \lean{⟨hp, hq⟩} formado por una
                demostración de cada componente, que se recuperan con
                \lean{.1} y \lean{.2}.
          \item El bicondicional \lean{p ↔ q} es un par de implicaciones,
                cuyas componentes se extraen con \lean{.mp} (de izquierda a
                derecha) y \lean{.mpr} (de derecha a izquierda).
          \item El existencial \lean{∃ x : A, P x} se demuestra con un par
                \lean{⟨a, ha⟩}: un testigo \lean{a} y una
                demostración \lean{ha : P a}. Es a la conjunción lo que el
                cuantificador universal es a la implicación: un par en el
                que el tipo de la segunda componente depende de la primera.
          \item La disyunción \lean{p ∨ q} es la excepción de este grupo: no
                es un par sino una suma, con dos constructores,
                \lean{Or.inl hp} y \lean{Or.inr hq}, según qué componente se
                demuestre. Construirla exige elegir un lado. Usarla, en
                cambio, obliga a tratar los dos casos posibles.
        \end{itemize}
\end{itemize}

Los corchetes angulares \lean{⟨ , ⟩} que acaban de aparecer son el
constructor anónimo: una notación uniforme para construir términos de
cualquier tipo con un solo constructor. Los tres primeros casos del segundo
grupo son instancias de lo mismo, un par, con una o dos componentes que
pueden depender entre sí, y por eso los tres se escriben igual; la
disyunción, con sus dos constructores, queda fuera. Esta notación aparecerá
constantemente, también para los elementos de las estructuras que definamos.

\subsection{Tipos inductivos}\label{sec:tipos-inductivos}

La tercera idea es el mecanismo con el que se crean tipos nuevos: los
tipos inductivos. Un tipo inductivo se define dando una lista de
constructores, y sus términos son exactamente los que se obtienen
aplicando constructores un número finito de veces. El ejemplo clásico son
los números naturales, definidos en la biblioteca de Lean esencialmente así:

\begin{lstlisting}
inductive Nat where
  | zero : Nat              -- el cero es un natural
  | succ : Nat → Nat       -- el sucesor de un natural es un natural
\end{lstlisting}

Cada línea que empieza por \lean{|} declara un constructor con su tipo. La
definición debe leerse como: los términos de tipo \lean{Nat} son
\lean{Nat.zero}, \lean{Nat.succ Nat.zero}, \lean{Nat.succ (Nat.succ
Nat.zero)}, etcétera, y nada más. De este ``y nada más'' Lean
extrae automáticamente dos herramientas:

\begin{itemize}
  \item Un \emph{principio de recursión}: para definir una función sobre un
        tipo inductivo basta decir qué hace con cada constructor. Lo
        usaremos, por ejemplo, para definir la relación de satisfacción por
        recursión sobre la estructura de las fórmulas.
  \item Un \emph{principio de inducción}: para demostrar que todos los
        términos del tipo cumplen una propiedad basta demostrarla para cada
        constructor, suponiéndola para sus argumentos (hipótesis de
        inducción). La inducción usual sobre $\mathbb{N}$ es el caso
        particular para \lean{Nat}; nosotros usaremos, además, la
        inducción estructural sobre fórmulas.
\end{itemize}

Hay un uso de los tipos inductivos menos evidente y absolutamente central en
este trabajo: se pueden definir proposiciones inductivas, es decir,
familias de términos de \lean{Prop} generadas por constructores. Nuestra
noción de demostrabilidad ($\proves \varphi$) será una proposición inductiva
cuyos constructores son, literalmente, los axiomas y reglas del sistema
$\K$: una demostración en $\K$ queda representada como un árbol finito de
aplicaciones de constructores. La consecuencia es muy potente: el principio
de inducción asociado permite hacer inducción sobre las
demostraciones, que es exactamente la técnica con la que se prueba el
teorema de corrección (``si los axiomas son válidos y las reglas preservan
la validez, todo teorema es válido'').

Por último, las estructuras (\lean{structure}) son el caso particular
de tipo inductivo con un único constructor: sirven para empaquetar varios
campos en un solo objeto, como un marco de Kripke (un tipo de mundos y una
relación) o un modelo (un marco y una valoración). Si \lean{F} es una
estructura con un campo \lean{R}, se accede a él con la notación
\lean{F.R}; y una estructura puede extender a otra
(\lean{extends}), heredando sus campos, como hará \lean{KripkeModel} con
\lean{KripkeFrame}. Esto no es de especial relevancia, y se construirá con ánimo meramente ilustrativo.

\section{Definiciones y enunciados}\label{sec:escribir-lean}

Con el fundamento anterior, escribir Lean consiste en declarar términos y
sus tipos. Las declaraciones que usamos son cuatro:

\begin{itemize}
  \item \lean{def} introduce una definición con nombre. Por ejemplo,
        \lean{def EsPar (n : Nat) : Prop := ∃ k : Nat, n = 2 * k} define la
        propiedad de ser par: recibe un natural \lean{n} y devuelve la
        proposición que afirma que \lean{n} es el doble de algún
        \lean{k}. Entre el nombre y los dos puntos van los argumentos con su
        tipo; lo que sigue a los dos puntos es el tipo del resultado (aquí
        \lean{Prop}, porque lo definido es un enunciado y no un número), y
        lo que sigue a \lean{:=} es el término definido. Así se definirán
        propiedades en el capítulo \ref{sec:formalizacion} con nociones como la validez o la consistencia.
  \item \lean{theorem} declara un teorema: su ``tipo'' es el enunciado (una
        proposición) y su ``definición'' es la demostración (un término de
        ese tipo). Que teoremas y definiciones tengan la misma forma es un ejemplo más de correspondencia de Curry--Howard.
  \item \lean{example} es un teorema sin nombre: se verifica y se olvida.
        Lo usamos para ejemplos o verificaciones que no se necesitan después.
  \item \lean{variable} declara parámetros comunes a varias declaraciones.
        En nuestro código, por ejemplo, \lean{variable \{VP : Type\}} fija de una vez el
        tipo de las variables proposicionales para todo el desarrollo.
\end{itemize}

\vspace{3mm}

\underline{Argumentos implícitos}

En algunas declaraciones aparecen argumentos
entre llaves en lugar de entre paréntesis. Las llaves indican que el
argumento es implícito: quien usa la definición no lo escribe, y Lean
lo infiere del contexto. Por ejemplo, un lema sobre listas que afirme que
invertir dos veces devuelve la lista de partida,
\lean{theorem rev_rev \{α : Type\} (l : List α) : (l.reverse).reverse = l},
tiene el tipo \lean{α} implícito: al aplicarlo a una lista concreta de
naturales, Lean deduce por sí solo que \lean{α} es \lean{Nat} y escribirlo
sería redundante. En el código, la regla práctica que se ha seguido es marcar como
implícito todo argumento que pueda deducirse del tipo de los demás, y dejar
explícitos los que no. En el capítulo \ref{sec:formalizacion}, el tipo de variables proposicionales sobre el que se construye todo el desarrollo será un argumento implícito de casi todas las definiciones y teoremas.

\vspace{3mm}

\underline{Notación}

Lean permite introducir notación infija y prefija con
niveles de precedencia. Por ejemplo,

\begin{lstlisting}
local infixr:60 " ⟶ " => ModalFormula.imp
local prefix:70 "□"     => ModalFormula.box
\end{lstlisting}

declara que \lean{φ ⟶ ψ} abrevia \lean{ModalFormula.imp φ ψ} y que
\lean{□φ} abrevia \lean{ModalFormula.box φ}. El número tras los dos puntos
es la precedencia. A mayor número, más fuerte liga: así
$\nec\varphi \longrightarrow \psi$ se lee $(\nec\varphi) \longrightarrow
\psi$. La \lean{r} de \lean{infixr} indica asociatividad a la derecha (la
convención habitual para la implicación) y el modificador \lean{local}
restringe la notación al fichero actual. Como se discute en la sección
\ref{sec:form-sintaxis}, elegimos deliberadamente símbolos distintos de los
de Lean ($\longrightarrow$ frente a $\to$, $\sim$ frente a $\lnot$,
$\curlywedge$ frente a $\land$\dots) para que siempre sea visible si una
expresión es una fórmula modal o una proposición de Lean.

\vspace{3mm}

\underline{Definiciones por recursión}

Como anticipamos, una función sobre un
tipo inductivo puede definirse por casos sobre los constructores, mediante
ajuste de patrones (\emph{pattern matching}):

\begin{lstlisting}
def suma : Nat → Nat → Nat
  | Nat.zero,   n => n
  | Nat.succ m, n => Nat.succ (suma m n)
\end{lstlisting}

Cada línea \lean{| patrón => valor} cubre un constructor; Lean comprueba que
los casos son exhaustivos y que las llamadas recursivas se hacen sobre
argumentos estrictamente menores, lo que garantiza que la definición
termina. Sobre estas definiciones se apoya la noción de \emph{igualdad
definicional}: dos expresiones son definicionalmente iguales si una se
obtiene de la otra desplegando definiciones y computando. Por ejemplo,
\lean{suma Nat.zero n} es definicionalmente \lean{n}, sin más que aplicar la
primera cláusula. Las igualdades definicionales no necesitan demostración:
la táctica \lean{rfl} (sección \ref{sec:tacticas}) las cierra de inmediato,
y en la formalización señalaremos varios lemas cuya demostración completa
es, precisamente, \lean{rfl}. La relación de satisfacción del capítulo
\ref{sec:formalizacion} se define exactamente así, con una cláusula por cada
conectiva del lenguaje modal.

\section{El modo táctico y el InfoView}\label{sec:modo-tactico}

En principio, demostrar un teorema es escribir un término de su tipo, y para
enunciados sencillos puede hacerse directamente. Por ejemplo, el modus
ponens:

\begin{lstlisting}
theorem mp {p q : Prop} (h : p → q) (hp : p) : q := h hp
\end{lstlisting}

La demostración es la aplicación \lean{h hp}, tal como vimos en la sección
\ref{sec:tipos}: la implicación es una función y aplicarla es razonar.
Pero construir a mano el término de una demostración larga es impracticable.
Para eso existe el \emph{modo táctico}: escribiendo \lean{by} tras el
enunciado, en lugar de un término se da una sucesión de tácticas,
instrucciones que van construyendo el término por nosotros.

El modo táctico se entiende mejor con la noción de estado de la
demostración, que el InfoView muestra en todo momento. Un estado consta de
uno o varios objetivos (\emph{goals}); cada objetivo tiene un
contexto de hipótesis con nombre (\lean{n : Nat}, \lean{h : p},
\dots) y una meta, la proposición que falta demostrar, señalada con
el símbolo $\vdash$. Cada táctica transforma el estado: puede simplificar la
meta, añadir hipótesis, dividir el objetivo en varios o cerrarlo. La
demostración acaba cuando no quedan objetivos y el InfoView muestra el
mensaje ``\emph{Goals accomplished!}''.

Veámoslo con un ejemplo elemental: si $p$ es cierta, entonces $p$ se sigue
de cualquier otra cosa. En los comentarios se indica el estado de la
demostración antes y después de cada táctica, tal como lo mostraría el
InfoView.

\begin{lstlisting}
example (p q : Prop) : p → (q → p) := by
  -- ⊢ p → q → p
  intro hp
  -- hp : p
  -- ⊢ q → p
  intro hq
  -- hp : p,  hq : q
  -- ⊢ p
  exact hp
  -- Goals accomplished!
\end{lstlisting}

La meta inicial es una implicación anidada. La táctica \lean{intro hp}
supone el antecedente ---es el ``supongamos $p$'' de una demostración
informal--- y lo añade al contexto con el nombre \lean{hp}, dejando como
meta el consecuente; \lean{intro hq} hace lo propio con la segunda
implicación. La meta es entonces $p$, que es exactamente la hipótesis
\lean{hp}, y \lean{exact hp} la señala y cierra el objetivo. Estas dos
tácticas, junto con el resto de las empleadas en el trabajo, se detallan en
la sección \ref{sec:tacticas}.

Merece la pena subrayar la relación entre los dos modos: las tácticas no son
una lógica aparte, sino un lenguaje de órdenes para construir el mismo
término que podríamos escribir a mano. En el ejemplo, el término construido
es \lean{fun hp hq => hp}. A lo largo del trabajo combinamos ambos estilos
según cuál resulte más legible ---alguna demostración corta se escribe
directamente como término, igual que el modus ponens de arriba---, con
preferencia por el modo táctico con las hipótesis intermedias explícitas
(\lean{have}), que es el que más se parece a una demostración escrita en
prosa.

\section{Catálogo de tácticas empleadas}\label{sec:tacticas}

Esta sección recoge todas las tácticas que aparecen en la
formalización, agrupadas por función, con su significado matemático y una
referencia a algún lugar del capítulo \ref{sec:formalizacion} donde se
utilicen. La tabla \ref{tab:tacticas} sirve de resumen y de índice.

\begin{table}[htb]
\centering\small
\begin{tabular}{@{}lll@{}}
\toprule
Táctica & Lectura matemática & Uso destacado \\
\midrule
\lean{intro}       & sea\ldots / supongamos\ldots & en casi todos los lemas \\
\lean{exact}       & que es exactamente\ldots & en casi todos los lemas \\
\lean{apply}       & basta ver\ldots (razonar hacia atrás) & lema de verdad \\
\lean{have}        & afirmación intermedia & regla RE (§\ref{sec:form-sistema}) \\
\lean{rfl}         & por definición & \lean{satisfies_neg} (§\ref{sec:form-semantica}) \\
\lean{constructor} & separar $\leftrightarrow$, $\wedge$ en partes & \lean{satisfies_dia} \\
\lean{rintro}, \lean{rcases}, \lean{obtain} & descomponer hipótesis & Lindenbaum (§\ref{sec:form-completitud}) \\
\lean{refine}      & esquema con huecos & lema de existencia \\
\lean{by_contra}   & reducción al absurdo & completitud \\
\lean{by_cases}    & distinción de casos clásica & propiedades de los MCS \\
\lean{contradiction} & hipótesis contradictorias & tautologías del $\leftrightarrow$ \\
\lean{induction ... with} & inducción (estructural o sobre derivaciones) & corrección, lema de verdad \\
\lean{subst}       & sustituir usando una igualdad & teorema de deducción \\
\lean{simp}        & simplificación automática & un único uso (§\ref{sec:form-completitud}) \\
\bottomrule
\end{tabular}
\caption{Resumen de las tácticas empleadas en la formalización.}
\label{tab:tacticas}
\end{table}

\subsection{Tácticas básicas de manipulación de objetivos}

\textbf{\lean{intro}.} Si la meta es una implicación \lean{p → q} o un
universal \lean{∀ x, P x}, la táctica \lean{intro h} mueve el antecedente al
contexto como nueva hipótesis \lean{h} (o introduce un \lean{x} genérico) y
deja como meta el consecuente. Corresponde al ``supongamos que'' y al ``sea
$x$ arbitrario'' del lenguaje matemático. Admite varios nombres seguidos:
ante la meta \lean{⊢ p → q → p}, la táctica \lean{intro hp hq} deja el
contexto con \lean{hp : p} y \lean{hq : q} y la meta \lean{⊢ p}. Es, con
diferencia, la táctica más usada del trabajo: casi todas nuestras nociones
(tautología, validez, consistencia) son cadenas de cuantificadores e
implicaciones.

\textbf{\lean{exact}.} Cierra la meta proporcionando un término que la
demuestra, construido con las hipótesis disponibles. En el caso más simple
el término es una hipótesis: con \lean{hp : p} y meta \lean{⊢ p}, basta
\lean{exact hp}. Pero puede ser cualquier término del tipo adecuado: con
\lean{h : p → q} y \lean{hp : p} ante la meta \lean{⊢ q}, se cierra con
\lean{exact h hp}, aplicando la implicación a la hipótesis. En el trabajo
aparece a menudo en esta segunda forma, encadenando varias aplicaciones en
un solo término.

\textbf{\lean{apply}.} Razona hacia atrás: si \lean{h : p → q} y la meta es
\lean{q}, la táctica \lean{apply h} reduce la meta a \lean{p} (``para ver
$q$, por $h$ basta ver $p$''). Es la contrapartida de \lean{exact}: en lugar
de dar el término completo de una vez, se avanza dejando pendiente lo que
falta. La usamos, entre otros sitios, en el caso de la implicación del lema
de verdad (sección \ref{sec:form-completitud}).

\textbf{\lean{have}.} Introduce una afirmación intermedia con su propia
demostración: \lean{have h : p := término} o bien
\lean{have h : p := by ...}. Tras ella, \lean{h} queda disponible en el
contexto como cualquier otra hipótesis. Es la táctica que estructura las
demostraciones largas como sucesiones de pasos con nombre(``primero observamos\dots, de donde\dots''). La
demostración de la regla RE en la sección \ref{sec:form-sistema} es el
ejemplo más evidente: siete \lean{have} encadenados que se leen exactamente igual
que la demostración en papel.

\textbf{\lean{rfl}.} Cierra metas de la forma \lean{a = b} o \lean{a ↔ b}
cuando ambos lados son definicionalmente iguales (sección
\ref{sec:escribir-lean}): por ejemplo, la meta \lean{⊢ 2 + 2 = 4}, que se
resuelve calculando. Que un lema se demuestre con \lean{rfl} aporta además información
importante: significa que no hay contenido matemático nuevo, solo se despliegan
definiciones. Lo veremos en \lean{satisfies_neg} y en varios casos de la
inducción de \lean{cpl_eval_iff_satisfies}, en la sección
\ref{sec:form-semantica}.

\subsection{Construir y descomponer estructuras lógicas}

\textbf{\lean{constructor}.} Cuando la meta es una conjunción o un
bicondicional, la divide en sus componentes, generando un objetivo por cada
una: ante \lean{⊢ p ∧ q} deja los dos objetivos \lean{⊢ p} y \lean{⊢ q}, y
ante \lean{⊢ p ↔ q}, las dos implicaciones. El InfoView los muestra como
objetivos separados, para delimitar la
demostración de cada uno. Alternativamente puede darse el término par
directamente con el constructor anónimo (sección \ref{sec:tipos}): la
demostración final del teorema de adecuación es, simplemente,
\lean{⟨soundness, completeness⟩}.

\textbf{\lean{rcases}, \lean{rintro}, \lean{obtain}.} Son las tácticas
duales de \lean{constructor}: en lugar de dividir la meta, descomponen
hipótesis compuestas mediante patrones. Con
\lean{h : ∃ x, P x ∧ Q x}, la táctica \lean{rcases h with ⟨x, hP, hQ⟩}
sustituye \lean{h} por el testigo \lean{x} y las dos componentes de la
conjunción, \lean{hP : P x} y \lean{hQ : Q x}. Los patrones pueden anidarse
tanto como haga falta. Con \lean{h : p ∨ q}, en cambio,
\lean{rcases h with hp | hq} genera dos objetivos, uno por cada caso: la
barra separa alternativas y los corchetes angulares agrupan componentes
simultáneas. La táctica \lean{rintro} combina \lean{intro} con esta
descomposición en un solo paso, y \lean{obtain} hace lo propio con el
resultado de aplicar un teorema: \lean{obtain ⟨x, hP, hQ⟩ := teorema hip}
aplica el teorema y descompone lo que devuelve. La lectura matemática es
siempre la misma: ``sea $x$ el objeto que nos proporciona el teorema, con
tales propiedades''. En el trabajo aparecen sobre todo al usar el lema de
Lindenbaum y el lema de existencia (sección \ref{sec:form-completitud}).

\textbf{\lean{refine}.} Es un \lean{exact} con huecos: se escribe el
esqueleto del término y se dejan \lean{?_} en las partes pendientes, que se
convierten en nuevos objetivos. Ante la meta \lean{⊢ p ∧ q}, teniendo ya
\lean{hp : p}, la táctica \lean{refine ⟨hp, ?_⟩} deja como único objetivo
\lean{⊢ q}. Resulta ideal cuando la parte creativa, elegir el testigo de
un existencial, es previa a las comprobaciones rutinarias, que es
exactamente la situación del lema de existencia de la sección
\ref{sec:form-completitud}.

\subsection{Razonamiento clásico}

\textbf{\lean{by_contra}.} Reducción al absurdo: para demostrar \lean{p}, la
táctica \lean{by_contra h} añade la hipótesis \lean{h : ¬p} y cambia la meta
a \lean{False}, de modo que basta llegar a una contradicción. Es un
principio genuinamente clásico que equivale a la eliminación de la doble
negación. Su papel en este trabajo es muy relevante: la propia
demostración de completitud es una gran reducción al absurdo (sección
\ref{sec:form-completitud}), y en pequeño aparece cada vez que hay que pasar
de una negación a una afirmación, como en la inferencia
$\lnot\forall x\, \lnot P(x) \Rightarrow \exists x\, P(x)$, que no es
demostrable intuicionistamente. Volvemos sobre esta cuestión en la sección
\ref{sec:mathlib}.

\textbf{\lean{by_cases}.} Distinción de casos sobre una proposición
cualquiera: \lean{by_cases h : p} abre dos objetivos, uno con \lean{h : p} y
otro con \lean{h : ¬p}. Es la herramienta natural para razonar por casos sin
disponer de una disyunción previa: en el trabajo la usamos para las
propiedades de los conjuntos maximales consistentes (sección
\ref{sec:form-completitud}), del tipo ``o bien $\Gamma$ deduce $\varphi$, o
bien no lo hace''.

\textbf{\lean{contradiction}.} Cierra el objetivo cuando el contexto
contiene una contradicción explícita: por ejemplo, \lean{h : p} junto con
\lean{h' : ¬p}, o una hipótesis \lean{h : False}. En tal caso la meta, sea
cual sea, queda demostrada.

Relacionadas con lo anterior están dos funciones (no tácticas) que usamos
para rematar razonamientos: \lean{False.elim h}, que a partir de
\lean{h : False} demuestra cualquier cosa (\emph{ex falso quodlibet}), y
\lean{absurd h h'}, que hace lo mismo a partir de \lean{h : p} y
\lean{h' : ¬p}.

\subsection{Inducción y reescritura}

\textbf{\lean{induction ... with}.} Aplica el principio de inducción del
tipo inductivo correspondiente, abriendo un objetivo por constructor; la
sintaxis \lean{with | caso₁ ... | caso₂ ...} permite nombrar en cada caso
los argumentos y las hipótesis de inducción. Un ejemplo con listas:

\begin{lstlisting}
def longitud {α : Type} : List α → Nat
  | []      => 0
  | _ :: xs => longitud xs + 1

example {α : Type} (l : List α) : longitud (l ++ []) = longitud l := by
  induction l with
  | nil           => rfl
  | cons x xs ih  => rw [List.cons_append, longitud, longitud, ih]
\end{lstlisting}

El tipo \lean{List α} tiene dos constructores: la lista vacía \lean{[]} y la
que antepone un elemento a otra lista, \lean{x :: xs}. La inducción abre por
tanto dos objetivos: el caso \lean{nil}, que es la base y aquí se cierra de forma inmediata con \lean{rfl}, y el caso \lean{cons}, que recibe el elemento \lean{x}, la cola
\lean{xs} y la hipótesis de inducción \lean{ih}, que afirma la propiedad
para \lean{xs}. Lo mismo ocurre con cualquier otro tipo inductivo: hay un
caso por constructor, y los constructores recursivos vienen acompañados de
su hipótesis de inducción.

\begin{itemize}
  \item \emph{Inducción estructural sobre fórmulas} (tipo
        \lean{ModalFormula}): los casos son los cuatro constructores
        ---átomo, $\bot$, implicación y caja---, y los dos últimos traen
        hipótesis de inducción sobre las subfórmulas. Así se demuestran el
        lema puente de la corrección y el lema de verdad (secciones
        \ref{sec:form-correccion} y \ref{sec:form-completitud}).
  \item \emph{Inducción sobre derivaciones} (tipos \lean{Provable} y
        \lean{DeducibleFrom}): como una demostración formal es también un
        objeto inductivo, se puede razonar ``por inducción sobre la
        demostración'', con un caso por axioma o regla. De especial relevancia para la corrección, el teorema de
        deducción y la monotonía de la deducibilidad.
  \item Con el modificador \lean{generalizing}: cuando la propiedad que se
        demuestra por inducción depende de otra variable que cambia dentro
        de los casos, esa variable debe generalizarse antes de inducir, o
        las hipótesis de inducción quedarán fijadas a su valor actual y no
        servirán. Escribir \lean{induction x generalizing y} lo consigue. Es
        el análogo formal de la precaución, habitual en papel, de enunciar
        la hipótesis de inducción ``para todo $y$''. En el lema de verdad
        resulta imprescindible, y allí explicamos por qué (sección
        \ref{sec:form-completitud}).
\end{itemize}

\textbf{\lean{subst}.} Dada una hipótesis de igualdad \lean{heq : x = a},
elimina \lean{x} sustituyéndolo por \lean{a} en la meta y en las demás
hipótesis, y descarta la igualdad. Es lo que en papel se hace al decir
``pero entonces $x$ es $a$'', y se usa cuando una distinción de casos
ha dejado como hipótesis la identidad de dos objetos que se trataban por
separado.

\textbf{\lean{simp}.} Es la principal táctica de automatización de Lean:
reescribe la meta (o una hipótesis, con \lean{simp at h}) usando una base de
ecuaciones marcadas como reglas de simplificación, y cierra el objetivo si
llega a algo trivial. Es muy cómoda pero tiene un precio: oculta los pasos que da,
de modo que su uso indiscriminado produce demostraciones cortas pero
ilegibles. Por eso en este trabajo, aparece una única vez y en un papel muy concreto: para descartar una hipótesis imposible, la pertenencia de una
fórmula al conjunto vacío (sección \ref{sec:form-completitud}). Todo lo
demás está demostrado de forma explícita.

\section{Mathlib, la lógica clásica y el lema de Zorn}\label{sec:mathlib}

Nuestro fichero comienza con dos importaciones:

\begin{lstlisting}
import Mathlib.Tactic
import Mathlib.Order.Zorn
\end{lstlisting}

La primera carga el modo táctico extendido de Mathlib; la segunda, el lema
de Zorn. De Mathlib usamos tres cosas, que repasamos a continuación.

\vspace{3mm}

\underline{Conjuntos}

El tipo \lean{Set α} representa los subconjuntos de un
tipo \lean{α}; técnicamente, un conjunto es su función
característica, \lean{Set α := α → Prop}, y \lean{x ∈ A} es la proposición
\lean{A x}. Sobre él están definidos \lean{⊆}, \lean{∪}, \lean{∅}, la
inserción de un elemento \lean{insert a A} (que usamos constantemente para
``añadir una fórmula a un conjunto de hipótesis'') y la unión de una familia
\lean{⋃₀ C} (unión de todos los conjuntos que pertenecen a la colección
\lean{C}). De la biblioteca tomamos lemas elementales de manipulación cuyo
nombre describe su contenido: \lean{Set.mem_insert} ($a \in
\mathrm{insert}\ a\ A$), \lean{Set.mem_insert_iff} ($x \in \mathrm{insert}\
a\ A \leftrightarrow x = a \lor x \in A$), \lean{Set.subset_insert},
\lean{Set.mem_insert_of_mem}, \lean{Set.mem_sUnion} y
\lean{Set.subset_sUnion_of_mem}. Aprender a encontrar estos lemas
(por su nomenclatura sistemática, con la búsqueda del editor o con
herramientas como \lean{exact?}) ha sido parte del aprendizaje de Lean.

\vspace{3mm}

\underline{Lógica clásica}

El núcleo de Lean es constructivo: por defecto no
se dispone del tercero excluido. Mathlib, como la práctica matemática
habitual, trabaja en lógica clásica, axiomatizada mediante
\lean{Classical.choice}; de ahí derivan el tercero excluido y las tácticas
\lean{by_contra} y \lean{by_cases} que acabamos de describir. Nuestro
desarrollo es clásico sin complejos, igual que el texto de referencia
\cite{blackburn}: la propia semántica evalúa cada fórmula a verdadera o
falsa, la eliminación de la doble negación es una de nuestras tautologías
básicas, y la completitud usa el lema de Zorn, equivalente al axioma de
elección. No es una limitación de Lean, sino una elección: existe una
lógica modal intuicionista, pero no es el objeto de este trabajo.

\vspace{3mm}

\underline{El lema de Zorn}

Del orden por inclusión entre conjuntos, Mathlib
ofrece la siguiente versión del lema de Zorn, que es la que aplicaremos en
el lema de Lindenbaum (sección \ref{sec:form-completitud}):

\begin{lstlisting}
theorem zorn_subset_nonempty (S : Set (Set α))
    (H : ∀ c ⊆ S, IsChain (· ⊆ ·) c → c.Nonempty →
         ∃ ub ∈ S, ∀ s ∈ c, s ⊆ ub)
    (x) (hx : x ∈ S) :
    ∃ m, x ⊆ m ∧ Maximal (· ∈ S) m
\end{lstlisting}

Leído con calma: \lean{S} es una familia de conjuntos. La hipótesis \lean{H}
pide que toda cadena no vacía \lean{c} contenida en \lean{S}, (esto
es, toda subfamilia totalmente ordenada por $\subseteq$, que es lo que
expresa \lean{IsChain (· ⊆ ·) c}, donde \lean{(· ⊆ ·)} es la notación
anónima de Lean para la relación de inclusión) tenga una cota superior
dentro de \lean{S}, La conclusión da, para cada \lean{x ∈ S}, un
elemento \lean{m} de la familia con \lean{x ⊆ m} que es maximal en ella.

El argumento \lean{x} se declara sin indicar sutipo, \lean{(x)}: Lean lo deduce de la hipótesis \lean{hx : x ∈ S}, que obliga a que sea un \lean{Set α}. A su vez, la maximalidad se expresa con el predicado \lean{Maximal (· ∈ S) m}, que significa que \lean{m} cumple la propiedad de pertenecer a \lean{S} y que ningún elemento que la cumpla lo contiene estrictamente, De él extraeremos las dos piezas que necesitamos con
\lean{.1} y \lean{.2}, la pertenencia y la maximalidad propiamente dicha.

En nuestro caso \lean{S} será la familia de los conjuntos
consistentes de fórmulas y la cota superior de una cadena, su unión. El trabajo matemático consistirá en demostrar que esa unión sigue siendo
consistente.
